% Generated draft: requires mathematical and source review.
\documentclass[11pt]{article}
\usepackage[margin=25mm]{geometry}
\usepackage{fontspec}
\setmainfont{Latin Modern Roman}
\usepackage{amsmath,amssymb,mathtools,unicode-math}
\setmathfont{Latin Modern Math}
\usepackage{enumitem,needspace,xurl,hyperref}
\setlength{\parindent}{0pt}
\setlength{\parskip}{5pt}
\setlength{\emergencystretch}{3em}
\newcommand{\N}{\mathbb{N}}
\newcommand{\Z}{\mathbb{Z}}
\newcommand{\Q}{\mathbb{Q}}
\newcommand{\R}{\mathbb{R}}
\newcommand{\C}{\mathbb{C}}
\newcommand{\F}{\mathbb{F}}
\newcommand{\A}{\mathbb{A}}
\newcommand{\PP}{\mathbb{P}}
\newcommand{\E}{\mathbb{E}}
\newcommand{\eps}{\varepsilon}
\newcommand{\dd}{\,\mathrm{d}}
\begin{document}
\section*{UnsolvedMath ID: 30005188}
\textbf{Source problem code:} catalogue entry
\subsection{Definitions and mathematical statement}
Let \((A,H,D)\) be a locally compact spectral triple, with \(*\)-representation \(\pi:A\to\mathcal B(H)\) and self-adjoint operator \(D\). For a compact operator \(T\) on \(H\), let \(\mu_n(T)\), \(n\geq 1\), denote its singular values in nonincreasing order. For \(p>0\), define the weak Schatten class by
\[
\mathcal L_{p,\infty}(H)=\left\{T\text{ compact}:\sup_{n\geq1}n^{1/p}\mu_n(T)<\infty\right\}.
\]
Thus membership expresses the singular-value bound \(\mu_n(T)=O(n^{-1/p})\).

The source's dimension terminology appears to mean that, for a candidate parameter \(d>0\), a compact spectral triple is called \(d\)-dimensional when
\[
(1+D^2)^{-d/(2p)}\in\mathcal L_{p,\infty}(H)
\qquad\text{for every }p>0,
\]
and a locally compact spectral triple is called \(d\)-dimensional when
\[
\pi(a)(1+D^2)^{-d/(2p)}\in\mathcal L_{p,\infty}(H)
\qquad\text{for every }p>0\text{ and every }a\in A.
\]
The exponent \(-d/(2p)\) is a reconstruction of damaged mathematical typography in the supplied original statement and requires comparison with the report. These conditions are recorded as the source's terminology; the record does not specify a minimality or sharpness condition that would single out an exact dimension.

\textbf{Requested endpoint problem.} Establish critical Cwikel-type estimates at \(p=2\) for noncompact spectral triples, including estimates for the Euclidean operator
\[
T_f=M_f(1-\Delta)^{-d/4}.
\]
Here, in the Euclidean interpretation, \(d\) is a positive integer, \(H=L^2(\R^d)\), \(M_f\) denotes multiplication by a measurable function \(f\), and \(\Delta=\sum_{j=1}^d\partial_j^2\) is taken with the nonpositive self-adjoint convention. The power of \(1-\Delta\) is defined by spectral calculus. Initially, the product is defined on the natural domain
\[
\left\{u\in L^2(\R^d):f(1-\Delta)^{-d/4}u\in L^2(\R^d)\right\};
\]
its boundedness or compactness requires hypotheses on \(f\). The Euclidean space and Laplacian sign convention make the displayed expression explicit; the supplied record does not itself state them in this detail.

\textbf{Unresolved formulation.} The record does not specify the admissible class of noncompact spectral triples, the function space and norm for \(f\), the endpoint operator ideal or singular-value bound, or the allowed dependence of constants. Consequently, it does not determine a single quantified endpoint inequality. In particular, no estimate for every measurable \(f\), or for every locally compact spectral triple, is asserted here. The weak Schatten conditions above explain the dimension motivation but do not supply the missing endpoint hypotheses or conclusion.

The original discussion asks how to compute spectral-triple dimension, describes Cwikel-type estimates as a tool, and distinguishes the pre-critical range \(p>2\), the post-critical range \(p<2\), and the difficult critical case \(p=2\). It also connects the Euclidean operators to Cwikel--Lieb--Rozenblum inequalities for numbers of bound states, described there as eigenvalues outside the essential spectrum, of Schr\"odinger operators. No specific Schr\"odinger operator or counting inequality is supplied.
\subsection{Short English statement}
The requested problem concerns critical \(p=2\) Cwikel-type estimates used to study the dimension of noncompact spectral triples, especially estimates for \(M_f(1-\Delta)^{-d/4}\) in Euclidean space. The supplied record omits the function spaces, endpoint bound, and structural hypotheses needed for a precise theorem. Its open status and verification claims are supplied curation, not independently verified facts.
\subsection{Sources}
\begin{itemize}
\item Local UnsolvedMath registry, supplied complete source record: permanent ID \(30005188\), source problem code \texttt{catalogue entry}, title ``Critical Cwikel Estimates for Noncompact Spectral Triples.'' The record gives proposed year 2022 and no named proposer.
\item Supplied report citation: \emph{Non-Commutative Geometry and Cyclic Homology} (2023), Oberwolfach Reports, pp.~1987--2058, DOI \texttt{10.4171/owr/2022/35}. \url{https://doi.org/10.4171/owr/2022/35}. This is also the supplied literature source labelled ``original Oberwolfach report.'' The record attributes the problem to workshop open-problem discussions documented in Oberwolfach Reports.
\item Supplied literature source: \emph{Dimensions and singular traces for spectral triples}. \url{https://arxiv.org/abs/math/0202108}. Its precise support for this endpoint problem is not specified in the record.
\item The original excerpt attributes Cwikel-type theory to Cwikel [2], Birman and Solomyak [1], Weidl [4], and Levitina, Sukochev and Zanin [3]. The corresponding bibliography entries are absent from the supplied record; the description of [3] as ``recent'' belongs to the source discussion.
\item Literature assessment date: \texttt{2026-08-22}. This is supplied dated curation, not an independently verified assessment. It labels the problem open, states that spectral dimension can be computed from summability and zeta or singular-trace growth in many examples, describes the critical Cwikel-type case as unresolved, and states that there is no general computation theorem for arbitrary spectral triples. The supplied background expressly calls for independent verification of time-sensitive status and cited claims.
\item The registry labels the statement \texttt{corrected\_verified} and reports a check against the official Oberwolfach report. These are supplied provenance claims; no independent source consultation or mathematical verification is claimed in this draft.
\end{itemize}
\end{document}
